% -----------------------------------------------------------
\section{Introduction} \label{sec:intro}
% ------------------------------------------------------------

\cite{lookout} proposed the \textit{lookout} algorithm for anomaly detection using a kernel density estimate with bandwidth based on persistence homology. Anomalies are identified by fitting a Generalized Pareto Distribution (GPD) to the negative logarithm of the smallest leave-one-out kernel density estimates, and declaring points with low GPD probabilities as anomalies.

In this paper, we explore some of the theoretical properties of the lookout algorithm, particularly the choice of bandwidth based on the Rips death diameters, and the consistency of the kernel density estimates produced by this choice. This leads us to propose a modified lookout algorithm using a different bandwidth, although still based on the Rips death diameters. The modified algorithm improves the performance and robustness of the original algorithm, under a broader range of conditions.

First we introduce some background material, before describing both the original lookout algorithm, and our proposed modifications.

\subsection{Anomalies and surprisals}

There is no broadly accepted definition for an anomaly. \citet{hawkins1980identification} defines an anomaly or outlier as \textit{an observation which deviates so much from other observations as to arouse suspicion it was generated by a different mechanism.} \citet{Barnett1978} define an anomaly or outlier as \textit{an observation (or a subset of observations) which appears to be inconsistent with the remainder of that set of data.} Both these definitions agree that anomalies are rare and different from other observations. Usually that is applied by fitting some (conditional) probability distribution to the set of observations, and considering whether each observation can be considered a likely random draw from the probability distribution. We formalize this in the following definition.

\begin{defn}[Anomalies]
  Given a set of observations $\bm{y}_{i} \in \mathbb{R}^m$ for $i \in \{1, \dots, n\}$, and a probability density function $f$, we define observation $\bm{y}_i$ as anomalous if
  \begin{equation}
    \int \mathbbm{1}(f(\bm{u}) < f(\bm{y}_{i})) d\bm{u} < \alpha,
  \end{equation}
  for some small value of $\alpha > 0$, where $\mathbbm{1}$ is the indicator function.
\end{defn}

\begin{defn}[Surprisals]
  Given a probability density function $f$, the \emph{surprisal} of an observation $\bm{y}$ drawn from the density $f$ is defined as $-\log f(\bm{y})$.
\end{defn}
These are better known as  ``log scores'' in statistics \citep{Good1952-zw}, but we prefer the term ``surprisals'' from information theory \citep{Stone2022}, as it provides a measure of how surprising an observation is, given the probability density function $f$. Specifically, it takes large values for observations that are unlikely under $f$. 

Let $S = -\log f(\bm{Y})$ be the random variable representing the surprisal of a random vector $\bm{Y}$ drawn from $f$, and let $G(s) = \text{Pr}(S\le s)$ be its distribution function. Then the surprisal probability of observation $\bm{y}_i$ is given by
$$p_i = 1 - G(-\log f \bm{y}_i)$$
and an observation is an anomaly iff $p_i < \alpha$.

\subsection{Extreme Value Theory}

Consider $n$ independent and identically distributed random variables $X_1, \dots, X_n$ having a distribution function $F$, with their maximum $M_n = \max \{X_1, \dots, X_n\} $. Extreme Value Theory (EVT) analyses and characterizes the behavior of $M_n$. One of the fundamental results of EVT is the Fisher-Tippet-Gnedenko Theorem. We restate a version of it from \citet[][Theorem 3.1.1]{coles2001introduction}.

% \begin{equation}\label{eq:evt1}
%     M_n = \max \{X_1, \, \ldots, \, X_n\} \, .
% \end{equation}
% If $F$ is known, the distribution of $M_n$ can be derived for any value of $n$ as
% \begin{align}\label{eq:evt2}
%     P\{M_n \leq z \} & = P\{X_1 \leq z, \, \ldots, \, X_n \leq z \} \, ,  \\
%     & = P\{X_1 \leq z\} \times \ldots \times P\{X_n \leq z\}\, , \\
%     & =  \left(F(z)\right)^n \, .
% \end{align}
% However, $F$ is not known in practice. This gap is filled by Extreme Value Theory, which studies approximate families of models for $F^n$ so that extremes can be modeled and uncertainty quantified. Similar to the Central Limit Theorem for means, the Fisher-Tippet-Gnedenko Theorem, also known as the the extreme value theorem or the extremal types theorem states that under certain conditions, a scaled maximum $\frac{M_n - a_n}{b_n}$, if $a_n$ and $b_n$ exist, have certain limit distributions.

\begin{theorem}[Fisher-Tippett-Gnedenko] \label{thm:FisherTippett}
  If there exist sequences of constants $\{a_n>0\}$ and $\{b_n\}$ such that
  $$ P\left\{ (M_n - b_n)/a_n \leq z \right\} \rightarrow G(z) \quad \text{as} \quad n \to \infty,
  $$
  for a non-degenerate distribution function $G$, then $G$ is a member of the Generalized Extreme Value (GEV) family
  \begin{equation}\label{eq:EVT4}
    G(z) = \exp\left\{ -\left[ 1 + \xi\left(\frac{z - \mu}{\sigma} \right)\right]^{-1/\xi} \right\}\, ,
  \end{equation}
  where the domain of the function is $\{z: 1 + \xi (z - \mu)/\sigma >0 \}$ and the parameters $\mu, \xi \in \mathbb{R}$ and $\sigma > 0$.
\end{theorem}
The parameters $\mu$, $\sigma$ and $\xi$ are called the location, scale and shape parameters respectively. The shape parameter $\xi$ determines the behavior of the tails and differentiates between the Gumbel, Fréchet and Weibull distributions: a Fréchet distribution is obtained when $\xi>0$, a Weibull distribution is obtained when $\xi < 0$, and a Gumbel distribution is obtained in the limit as $\xi\rightarrow 0$. Importantly, when $F$ has a finite upper bound, then $\xi < 0$.

A second result from EVT that we will use is Pickands theorem \citep{Pickands1975}, here restated from \citep[][Theorem 4.1]{coles2001introduction}.

\begin{theorem}[Pickands] Let $X_1, \dots, X_n$  be a sequence of independent random variables with a common distribution function $F$, and let
  $M_n = \max \{X_1, \ldots, X_n \}$. Suppose $F$ satisfies Theorem~\ref{thm:FisherTippett}, so that for large $n$, $P\{ M_n \leq z \} \approx G(z)$,
  where
  $$
  G(z) = \exp\left\{ -\left[ 1 + \xi\left(\frac{z - \mu}{\sigma} \right)\right]^{-1/\xi} \right\}\, ,
  $$
  for some $\mu, \xi \in \mathbb{R}$ and $\sigma >0$. Then for large enough $u$,
  \begin{equation}\label{eq:POT4}
    H(x) = \text{Pr}(X \le u+x \mid X > u) = 1 - \left( 1 + \frac{\xi x}{\sigma_u} \right)^{-1/\xi}\, ,
  \end{equation}
  where the domain of $H$ is $\{x: x > 0 \text{~and~} (1 + \xi x)/\sigma_u >0\}$, where $\sigma_u = \sigma + \xi(u- \mu)$.
\end{theorem}
The family of distributions defined by equation~\eqref{eq:POT4} is called the \textbf{Generalized Pareto Distribution} (GPD). The shape parameter $\xi$ is the same in both GPD and GEV distributions. Furthermore, the GEV parameters can be recovered from the GPD formulation by estimating a GPD on all observations above some threshold.

\subsection{Kernel Density Estimation}

Given a set of observations $\bm{y}_{i} \in \mathbb{R}^m$ for $i \in \{1, \dots, n\}$, a kernel density estimate is defined as \citep{chacon2018multivariate}
\begin{equation}\label{eq:mkde}
  \hat{f}(\bm{y}) = \frac{1}{n} \sum_{i=1}^n |\bm{H}|^{-1/2} K\big(\bm{H}^{-1/2}(\bm{y} - \bm{y}_i)\big),
\end{equation}
where $K$ is a square-integrable spherically-symmetric function, bounded below by 0, with a finite second-order moment and unit integral, and $\bm{H}$ is a symmetric $m\times m$ positive-definite matrix.

\begin{defn}[Consistency] Let $Y_1,\dots,Y_n$ be a sequence of independent random variables from a continuous distribution with square integrable density function $f$. Then \eqref{eq:mkde} is a \emph{consistent} estimator of $f$ if
  $$
  \lim_{n\rightarrow0} \text{E}\left[\left(\hat{f}(\bm{u}) - f(\bm{u})\right)^2 du\right] = 0
  $$
\end{defn}
The following result is provided by \citet[][p31--37]{chacon2018multivariate}.
\begin{theorem}
  For $\hat{f}$ to be a consistent estimator, $\bm{H}$ must be a symmetric positive-definite matrix that satisfies the following conditions:
  \begin{itemize}
    \item All elements of $\bm{H}$ approach zero as $n\rightarrow\infty$; and
    \item $n^{-1}|\bm{H}|^{-1/2} \rightarrow 0$ as $n\rightarrow\infty$.
  \end{itemize}
\end{theorem}

We consider matrices of the form $\bm{H} = h_n \bm{I}_m$, where $h_n\in\R$ is a function of the sample $\{\bm{y}_1,\dots,\bm{y}_n\}$
\begin{defn}
  We say that $h_n$ is an \emph{admissible bandwidth} if $(h_n)^{2/m}\rightarrow 0$ and $nh_n\rightarrow \infty$ as $n\rightarrow\infty$ almost surely.
\end{defn}
Thus, an admissible bandwidth leads to a consistent kernel density estimator.

\subsection{The lookout algorithm for anomaly detection}

The \textit{lookout} algorithm of \cite{lookout} has four inputs: the dataset $\{\bm{y}_1,\dots,\bm{y}_n\}$, where each $\bm{y}_i \in \mathbb{R}^m$; a parameter $\beta \in (0,1)$ indicating the proportion of points to be used for fitting the Generalized Pareto Distribution (GPD); a parameter $\alpha \in (0,1)$ signifying the cut-off threshold for determining an anomaly; and a logical variable \textit{unitize} indicating whether to scale the data. We briefly outline the steps below.

\begin{algorithm}
  \caption{Lookout algorithm}\label{alg:lookout}
  \begin{algorithmic}[1]
    \Require Data $\{\bm{y}_1, \dots, \bm{y}_n\}$, proportion $\beta \in (0,1)$, threshold $\alpha \in (0,1)$, \emph{unitize} $\in \{\text{True}, \text{False}\}$.
    \State If \textit{unitize}, scale the data to lie in the unit cube in $\mathbb{R}^m$.
    \State Compute the persistence homology barcode of the (possibly scaled) data cloud for dimension zero using the Vietoris-Rips diameter \citep{ghrist2008barcodes}.
    \State From the barcode, obtain the ordered death diameters $\{d_i\}_{i = 1}^n$ for connected components, and their successive differences $\Delta d_i = d_{i+1} - d_i$.
    \State Let $d_* = d_{i_*}$ where $i_* = \text{arg\! max} \{\Delta d_i \}_{i=1}^n$; that is, $d_*$ is the diameter corresponding to the lower point of the maximum $\Delta d_i$.
    \State Compute the kernel density estimates at the observations: $f_i = \hat{f}(\bm{y}_i)$, $i=1,\dots,n$, where $\hat{f}$ is given by \eqref{eq:mkde}, $\bm{H} = h_n\bm{I}_m$, and $h_n = (d_*)^{2/m}$.
    \State Compute the leave-one-out kde values $f_{-i} = \frac{1}{n-1} \big( n f_i - \bm{H}^{-1/2} K(\bm{0}) \big)$ for $i=1,\dots,n$.
    \State Fit a Generalized Pareto Distribution (GPD) to the largest $1-\beta$ of the surprisals $\{-\log f_i\}_{i=1}^n$, obtaining estimates of the location, scale and shape parameters $(\mu, \sigma, \xi)$.
    \State Calculate the probability of each observation by applying the estimated GDP to the leave-one-out kde values: $p_i = (1-\beta)P(-\log f_{-i} \mid \hat{\mu}, \hat{\sigma}, \hat{\xi})$.
    \State If $p_i < \alpha$, declare $\bm{y}_i$ anomalous.
  \end{algorithmic}
\end{algorithm}

The lookout algorithm provides density-based anomaly detection, where anomalies are identified as points with low density in the space of the data. If there are no anomalies in the data, the algorithm should identify $\alpha$ of the data as false positive ``anomalies''.

While \citet{lookout} demonstrated excellent results in applying their algorithm to a variety of simulated and real datasets, we see several potential areas for improvement that we aim to address in this paper. First, elements of the algorithm are not robust to outliers, including the min-max scaling, and the use of the largest gap in Rips death diameters. Second, the diagonal bandwidth matrix $\bm{H}$ does not account for correlations between variables, which can lead to suboptimal performance if the data have strong correlations. Third, it is not obvious that this bandwidth matrix choice will lead to a consistent kernel density estimator, particularly for distributions with heavy tails. Finally, the GPD estimation made minimal assumptions about the distribution of the surprisal values, but we know that these are bounded above and below, and that fact can be used to improve the estimated GPD.

In the following section, we will explore the conditions under which the kernel density estimator is consistent for $h_n=d_*$, and the implications for the lookout algorithm. In particular, we will show that the conditions are not satisfied for distributions with heavy tails such as a power distribution, although it seems likely that they are satisfied for lighter-tailed distributions such as the Gaussian distribution. We will also show that an alternative choice of bandwidth, based on an upper quantile of the Rips death diameters, satisfies these conditions for a wide range of distributions.

The main mathematical approach for proving when a bandwidth is admissible uses the notion of crackle, introduced in \cite{Crackle}, which describes the persistent homology of noise. They proved asymptotic results for the persistent homology of the union of unit balls centered on iid samples of various distributions. There are qualitative differences between sets of samples taken from the Gaussian distribution (which do not ``crackle'') compared to those from heavy-tailed distributions. We will adapt the methods to this setting where we wish to shrink the radius to $0$, albeit slowly.

%This paper proved asymptotic results about the

\subsection{The modified lookout algorithm}

As a result of these theoretical considerations, we propose three modifications to the lookout algorithm, described in detail in \cref{sec:modified}.

\begin{enumerate}
  \item Rather than use a max-min scaling of the data, we standardize the data using a robust covariance estimator. That is, we scale and rotate the data to ensure that the covariance matrix of the transformed data is approximately the identity matrix. This ensures that the data have zero pairwise correlations, which leads to more efficient kernel density estimates. It also leads to each variable having approximately unit variance, which provides a more robust scaling than min-max scaling.
  \item Rather than using the lower point of the largest gap in Rips death diameters, we use an upper quantile of the Rips death diameters. This avoids the problem where anomalies affect the largest gap in the Rips diameters that cause a persistent homology death. It also ensures the consistency of the kernel density estimates under relatively weak conditions.
  \item The estimated GPD shape parameter is constrained to be non-positive, which improves the stability of the GPD fitting. This is justified by noting that the kernel density estimates are bounded, so the limiting distribution under the Fisher-Tippett-Gnedenko theorem is a Weibull distribution, which has a negative shape parameter, $\xi$.
\end{enumerate}

Note that we do not need to estimate the density of the underlying distribution directly. Instead, we can estimate the density of a transformation of the data, and identify anomalies in the transformed space. Hence, there is no need to undo the scaling or rotation.

\subsection{Outline of the paper}

The remainder of this paper is organized as follows. In \cref{sec:theory}, we explore the theoretical properties of the lookout algorithm, particularly the choice of bandwidth based on Rips death diameters, and the consistency of the kernel density estimates produced by this choice. In \cref{sec:modified}, we propose a modified lookout algorithm that improves performance and robustness under a broader range of conditions. We provide some empirical experiments, demonstrating the performance of the modified lookout algorithm on simulated and real datasets, first in contrast to the original lookout (\cref{sec:comparewitholderversion}), and then in contrast to other anomaly detection methods (\cref{sec:compareWithOtherMethods}). Finally, in \cref{sec:conclusion}, we summarize our results and discuss future work.
